\section{Wybrane diagramy do opisu zaimplementowanych metod}
\subsection{Diagram przepływu danych (DFD)}
Diagram przepływu danych służy do wizualnego przedstawiania przepływu informacji w systemie, ukazując, jak dane przechodzą między procesami. W kontekście modelowania systemów klasyfikacyjnych, DFD pozwala na zilustrowanie sekwencji przetwarzania danych od wejścia do wyniku, co ułatwia zrozumienie struktury i dynamiki procesów, takich jak analiza danych w detekcji dezinformacji.

\subsection{Diagram aktywności (Activity Diagram)}

Diagram aktywności służy do wizualnego przedstawiania przepływu działań i procesów w systemie, ukazując, jak kolejne kroki lub zadania są wykonywane w określonej sekwencji. W kontekście modelowania systemów klasyfikacyjnych, takich jak analiza danych w detekcji dezinformacji, diagram aktywności pozwala na zilustrowanie logiki i dynamiki procesów, w tym sekwencji operacji (np. wczytywanie danych, trenowanie modelu, przetwarzanie wyników) oraz decyzji (np. wybór między różnymi metodami, takimi jak embeddings lub TF-IDF). Dzięki temu ułatwia zrozumienie struktury algorytmu i jego przepływu, szczególnie w złożonych zadaniach, jak ensemble learning.

\subsection{Diagram blokowy (Block Diagram)}
Diagram blokowy służy do wizualnego przedstawienia sekwencji procesów w systemie, ukazując, jak poszczególne etapy są ze sobą powiązane. W kontekście modelowania systemów klasyfikacyjnych, diagram blokowy umożliwia zilustrowanie ciągu przetwarzania danych od wejścia do wyniku, co ułatwia zrozumienie organizacji i przepływu działań, takich jak analiza danych w detekcji dezinformacji.